\section{Intrinsic multiplier ideals of the original smooth scheme affine opens}

\begin{lemma}[Nonzero generic stalk gives the actual nonzero affine ideals]\label{lem:sc0-original-generic-stalk-affine-nonzero}
\lean{EnvelopeProperty.Char0.affineIdeal_ne_bot_of_generic_stalk_ne_bot}
\leanok
On the original integral scheme $X$, let $I$ be an actual scheme
ideal sheaf whose actual generic stalk is nonzero. For every original
nonempty affine open $U$, the actual original ideal $I(U)$ is nonzero.
\end{lemma}
\begin{proof}
The generic point belongs to the nonempty open $U$. The actual
generic stalk ideal is the extension of $I(U)$ along the actual
generic germ. If $I(U)$ were zero, this extension would be zero,
contrary to its stated nonzero generic stalk.
\end{proof}

\begin{definition}[The original affine intrinsic multiplier ideal]\label{def:sc0-original-scheme-multiplier-affine-ideal}
\lean{EnvelopeProperty.Char0.schemeMultiplierAffineIdeal}
\leanok
\uses{lem:sc0-original-generic-stalk-affine-nonzero,thm:sc0-scheme-sections-function-field-equiv,thm:sc0-scheme-sections-function-field-fraction}
For the original structure morphism $f:X\to\operatorname{Spec} k$,
original ideal sheaf $I$ with nonzero generic stalk, and rational $q$,
define its affine multiplier ideal on every original nonempty affine
$U$ as follows. Map the whole actual ideal $I(U)$ along the actual
generic-germ equivalence $e_U:\Gamma(X,U)\simeq C_U\subset K(X)$,
form the full all-divisorial-DVR multiplier ideal on that literal
source $C_U$, and take its comap along $e_U$. On the empty open,
take its unique unit ideal.
\end{definition}

\begin{theorem}[Exact original nonempty principal compatibility]\label{thm:sc0-original-scheme-multiplier-principal-nonempty}
\lean{EnvelopeProperty.Char0.map_schemeMultiplierAffineIdeal_basicOpen_nonempty}
\leanok
\uses{def:sc0-original-scheme-multiplier-affine-ideal,thm:sc0-smooth-original-field-subalgebra,thm:sc0-scheme-sections-function-field-basic-localization,thm:sc0-full-multiplier-away-localization}
Suppose the actual original $f$ is smooth in characteristic zero.
For every original nonempty affine $U$ and nonempty intrinsic basic
open $D(s)$, extension along the actual original scheme restriction
maps this whole affine multiplier ideal to its whole actual affine
multiplier ideal on $D(s)$.
\end{theorem}
\begin{proof}
Use the actual generic-germ equivalences $e_U$ and $e_{D(s)}$.
Their exact restriction square identifies the actual section
restriction with the literal inclusion $C_U\subset C_{D(s)}$.
Both actual fraction fields are the original $K(X)$, original
smoothness proves smoothness of $C_U$, and the actual basic-open
localization theorem proves that this literal inclusion is
localization at $e_U(s)$.

The original ideal-sheaf restriction axiom and the actual square
give equality between the extension of $e_U(I(U))$ to $C_{D(s)}$
and the actual ideal $e_{D(s)}(I(D(s)))$. The full multiplier
localization theorem therefore gives exact equality of their whole
actual multiplier ideals. Finally transport that equality back along
the same actual square. Algebraically, for ring equivalences
$e:R\simeq S$ and $d:T\simeq P$ and a commuting square
$d a=b e$, extension of $e^{-1}(J)$ along $a$ is
$d^{-1}(J P)$. This follows by expressing comap along an equivalence
as map along its inverse and composing the actual ring maps. Apply
this identity to the actual restriction square to obtain precisely
the claimed original section-ring equality.
\end{proof}

\begin{theorem}[The complete actual basic-open ideal-sheaf axiom]\label{thm:sc0-original-scheme-multiplier-principal}
\lean{EnvelopeProperty.Char0.map_schemeMultiplierAffineIdeal_basicOpen}
\leanok
\uses{thm:sc0-original-scheme-multiplier-principal-nonempty}
For every actual original affine $U$ and every original section $s$,
including empty basic opens, the actual affine multiplier ideals
satisfy the full original ideal-sheaf basic-open restriction axiom.
\end{theorem}
\begin{proof}
If $D(s)$ is nonempty, then so is $U$, and the preceding theorem
applies with these actual nonempty opens. If $D(s)$ is empty, its
actual section ring is a subsingleton by the sheaf axiom for the
empty open. Its ideals are therefore identical, so the asserted
equality holds there as well.
\end{proof}
